%==============================================================================
% Paper E -- skeleton (bullets, to be polished into prose)
% Target venue: systems / measurement (e.g., a virtualization or measurement
%   workshop). This is the METHODS / INSTRUMENTATION paper of the portfolio:
%   it establishes how per-snapshot timing was measured and supplies the
%   snapshot-interval inertness finding the other papers (A, B, F, Plan 05)
%   rely on.
%
% Build: pdflatex skeleton.tex   (needs IEEEtran; if unavailable, change the
%   \documentclass line below to \documentclass[11pt]{article} and delete the
%   \IEEEauthorblock* macros.)
%==============================================================================
\documentclass[conference]{IEEEtran}
\usepackage{amsmath,booktabs,siunitx,url}

\title{Measure the Measurer: Per-Snapshot Timing Instrumentation and the\\
Snapshot-Interval Inertness Finding for Whole-VM Memory Capture}

\author{\IEEEauthorblockN{TODO Author list}
\IEEEauthorblockA{TODO Affiliation\\
\texttt{TODO@email}}}

\begin{document}
\maketitle

%------------------------------------------------------------------
\begin{abstract}
% TODO: polish into 150-200 words of prose. Bullet stub:
\begin{itemize}
  \item Problem: a whole-VM memory-capture pipeline exposes a single
        ``how often to sample'' knob (the requested guest sampling interval,
        \texttt{intervalMsec}). It was unclear what that knob actually
        controls, and whether tuning it changes capture behavior at all.
  \item Method: we added additive, low-overhead per-snapshot timing
        instrumentation -- six host \texttt{CLOCK\_REALTIME} timestamps per
        snapshot (\texttt{t0}\dots\texttt{t5}) bracketing
        \texttt{suspend}/\texttt{pmemsave}/\texttt{resume} -- and an
        orchestrator that derives a per-run \texttt{producer\_stats} summary.
  \item Central finding: the requested guest interval (\SI{500}{\milli\second}
        in the production campaign) is essentially \emph{inert}. The host
        snapshot cycle, dominated by QEMU \texttt{pmemsave}, dwarfs the
        interval, so lowering the interval below the snapshot-cycle floor is
        a no-op for capture throughput.
  \item Consequence: ``how often should we sample?'' is reframed as ``how
        cheaply can we snapshot?'' -- i.e.\ how to shrink \texttt{pmemsave}.
        This is the methodological pivot that Paper A (APF signal), Paper B
        (gated interval tuning), Paper F (streaming capture system), and the
        Plan 05 snapshot-throughput proposal all build on.
  \item Honesty: instrumentation also exposed two timing-correctness bugs (a
        missing-\texttt{bc} sleep fallback; a silent dump-cleanup failure)
        whose discovery is itself part of the contribution -- without the six
        timestamps these would have silently mis-calibrated everything
        downstream.
\end{itemize}
\end{abstract}

%------------------------------------------------------------------
\section{Introduction}
% Proposed introduction (bullets):
\begin{itemize}
  \item Hypervisor-external whole-VM memory capture (Paper A's APF pipeline)
        works by repeatedly pausing the guest, dumping all of guest physical
        RAM to disk via the QEMU monitor \texttt{pmemsave} command, and
        resuming. A configured interval, \texttt{intervalMsec}, nominally
        controls how often this happens.
  \item Before any sampling-rate study can be trusted, one must know what the
        interval knob actually sets. We argue there are \textbf{three distinct
        time axes} the original design conflated:
    \begin{itemize}
      \item Axis A -- \emph{guest-running interval}: the wall-clock gap between
            snapshots as measured by the guest's own monotonic clock (the VM
            does not run during \texttt{pmemsave}). This is the analysis frame
            spacing $\Delta t_{\mathrm{frame}}$.
      \item Axis B -- \emph{host wall-clock per snapshot}: real host time to
            produce one frame (suspend + \texttt{pmemsave} + resume).
      \item Axis C/D -- \emph{experiment throughput}: total snapshots
            $\times$ Axis B.
    \end{itemize}
  \item \texttt{intervalMsec} sets Axis A, \emph{not} Axis B. The post-resume
        \texttt{sleep intervalMsec/1000} is the only point in the producer
        loop where guest time advances; everything else happens with the VM
        paused.
  \item \textbf{Gap}: the pipeline shipped with no per-snapshot timing
        instrumentation, so it was impossible to tell (a) whether
        \texttt{intervalMsec} was honored, (b) how the per-snapshot host cost
        decomposes, or (c) whether tuning the interval changes anything that
        matters for capture throughput.
  \item \textbf{This paper} (the measurement/methods paper of the portfolio)
        builds that instrumentation, validates the three-axis model on real
        runs, and reports the headline empirical result: the requested guest
        interval is inert relative to the snapshot-cycle floor.
  \item \textbf{Contributions} (expanded in Sec.~\ref{sec:contrib}):
    \begin{itemize}
      \item C1: additive per-snapshot timing instrumentation
            (six host timestamps) and the \texttt{producer\_stats} run
            summary derived from them.
      \item C2: empirical separation of the three time axes, confirming
            \texttt{intervalMsec} sets only the guest-running interval
            (Axis A).
      \item C3: the snapshot-interval \emph{inertness} finding -- the host
            snapshot cycle (median \SI{9.75}{\second}, dominated by a median
            \SI{7.06}{\second} \texttt{pmemsave}) dwarfs the
            \SI{500}{\milli\second} request, so interval tuning below the
            floor is a no-op -- and the reframing of the sampling question as
            a snapshot-cost question.
      \item C4: bug-discovery as validation -- two timing-correctness defects
            surfaced only because the instrumentation existed.
    \end{itemize}
\end{itemize}

%------------------------------------------------------------------
\section{Background and Related Work}
\begin{itemize}
  \item VM introspection and memory forensics \cite{garfinkel2003vmi}:
        external observation of guest memory; the capture primitive here is a
        full-RAM \texttt{pmemsave} per snapshot.
  \item Checkpoint / snapshot cost \cite{TODO_libvirt_pmemsave}: the QEMU
        monitor \texttt{pmemsave} command writes a flat RAW image of all guest
        physical memory while the domain is paused. TODO: cite a
        checkpoint/snapshot-cost or live-migration study to frame why a full
        \SI{1}{\gibi\byte} dump dominates the per-snapshot cost.
        % TODO: add a citation key (e.g. \cite{TODO_snapshot_cost}) to refs.bib
        % for snapshot/checkpoint cost; none exists yet.
  \item Working-set / dirty-page tracking \cite{TODO_working_set_estimation}:
        hypervisors already track dirtied pages for migration. Position the
        ``shrink the snapshot'' future direction (Plan 05) against
        incremental dirty-bit capture, which would replace the full-RAM dump.
  \item TODO: cite measurement-methodology work on separating logical sampling
        rate from physical acquisition cost (the Axis-A vs Axis-B distinction
        is the methodological core of this paper).
\end{itemize}

%------------------------------------------------------------------
\section{Timing Instrumentation}
\label{sec:method}
\begin{itemize}
  \item \textbf{Capture primitive and loop.} Per snapshot the producer
        (\texttt{capture\_producer\_qemu\_pmemsave.sh}) runs:
        backpressure check $\rightarrow$ \texttt{virsh suspend} $\rightarrow$
        poll \texttt{domstate==paused} $\rightarrow$
        \texttt{pmemsave(0, ramBytes, file)} (VM frozen, writes a flat
        \SI{1}{\gibi\byte} RAW image) $\rightarrow$ flush $\rightarrow$
        enqueue $\rightarrow$ \texttt{virsh resume} $\rightarrow$ poll
        \texttt{domstate==running} $\rightarrow$
        \texttt{sleep intervalMsec/1000}. Guest RAM is \SI{1}{\gibi\byte}
        ($262{,}144$ pages of \SI{4096}{\byte}); each \texttt{pmemsave} writes
        exactly \SI{1}{\gibi\byte}.
  \item \textbf{Six timestamps.} An additive patch records, per snapshot, six
        host \texttt{CLOCK\_REALTIME} timestamps (nanosecond resolution via
        \texttt{date +\%s.\%N}):
        \texttt{t0} (before suspend), \texttt{t1} (after \texttt{paused}),
        \texttt{t2} (before \texttt{pmemsave}), \texttt{t3} (after
        \texttt{pmemsave} returns), \texttt{t4} (before resume), \texttt{t5}
        (after \texttt{running}). With \texttt{TIMING\_JSONL\_PATH} set the
        producer writes one JSONL record per snapshot (and a \texttt{seq:-1}
        record per backpressure event); when unset, ordinary captures are
        byte-for-byte unaffected. No new dependencies (Bash + GNU
        \texttt{date}).
  \item \textbf{Derived per-snapshot quantities.} From the six timestamps plus
        the next snapshot's \texttt{t0}:
        \[
          \texttt{pmemsave\_sec} = t_3 - t_2,\quad
          \texttt{suspend\_sec} = t_1 - t_0,
        \]
        \[
          \texttt{resume\_sec} = t_5 - t_4,\quad
          \texttt{host\_cycle\_sec} = t_0^{(n{+}1)} - t_0,
        \]
        \[
          \texttt{guest\_run\_sec} = t_0^{(n{+}1)} - t_5 \;\; (\text{Axis A}).
        \]
        \texttt{host\_cycle\_sec} is Axis B; \texttt{guest\_run\_sec} is
        Axis A; their ratio gives the per-snapshot VM-pause fraction.
  \item \textbf{Run summary (\texttt{producer\_stats}).} The orchestrator
        (\texttt{run\_timing\_instrumentation\_experiment.py}; CLI flags
        \texttt{--ram-mb}, \texttt{--duration}, \texttt{--interval-ms},
        \texttt{--keep-dumps}, \texttt{--output-json}) parses the JSONL and
        emits one JSON file with fields:
        \texttt{mean\_pmemsave\_sec},
        \texttt{mean\_host\_snapshot\_cycle\_sec},
        \texttt{mean\_suspend\_sec}, \texttt{mean\_resume\_sec},
        \texttt{mean\_guest\_run\_interval\_sec},
        \texttt{estimated\_vm\_pause\_fraction},
        \texttt{backpressure\_events}, \texttt{queue\_max\_depth}.
        Any field that cannot be measured is recorded as \texttt{null}, never
        guessed.
  \item \textbf{Producer-only by design.} Experiment~1 runs the producer with
        no consumer / delta calculation / streaming, so the timestamps measure
        pure capture cost. (The contended numbers in Sec.~\ref{sec:results}
        come from runs where a consumer competed for disk.)
  \item TODO: figure -- annotated timeline of one snapshot cycle showing
        \texttt{t0}\dots\texttt{t5}, the paused window ($t_0\to t_5$), and the
        single guest-running window (the post-resume sleep).
\end{itemize}

%------------------------------------------------------------------
\section{Contributions in Detail}
\label{sec:contrib}
\begin{itemize}
  \item \textbf{C1 -- instrumentation}: six additive host timestamps and the
        \texttt{producer\_stats} summary, with a \texttt{--dry-run} mode for
        host-free validation of paths and JSON shape.
  \item \textbf{C2 -- three-axis separation}: empirical confirmation that
        \texttt{intervalMsec} controls Axis A (guest-running interval) and is
        independent of Axis B (host wall-clock per snapshot).
  \item \textbf{C3 -- inertness finding}: the requested
        \SI{500}{\milli\second} interval is dominated by the host snapshot
        cycle; tuning it below the floor changes nothing about capture
        throughput. Reframes the sampling-rate question as a snapshot-cost
        question (feeds Paper B and Plan 05).
  \item \textbf{C4 -- validation by bug discovery}: the instrumentation caught
        a missing-\texttt{bc} sleep fallback (\texttt{intervalMsec} silently
        ignored) and a silent dump-cleanup failure (disk pressure inflating
        \texttt{pmemsave}); both would otherwise have corrupted downstream
        calibration invisibly.
\end{itemize}

%------------------------------------------------------------------
\section{Evaluation}
\label{sec:results}
% Proposed results (bullets with REAL numbers from the v3 campaign + the
% timing experiments). All numbers below are from the verified ground-truth
% block; do not recompute.
\begin{itemize}
  \item \textbf{Campaign context.} The production capture campaign (v3) is
        132 cells = 11 workloads $\times$ 3 durations
        $\{120, 300, 600\}\,\si{\second}$ $\times$ 4 replicates, all at a
        single requested interval of \SI{500}{\milli\second}; 128 cells
        completed with \texttt{exit\_status=="ok"}. Timing statistics below
        are over those \texttt{ok} cells.
  \item \textbf{Axis A is honored.} The guest indeed runs in
        $\approx$\SI{0.52}{\second} bursts between snapshots, close to the
        \SI{500}{\milli\second} request. So the \emph{guest} interval is
        honored; the dominant cost lives entirely on the host side.
  \item \textbf{The host snapshot cycle dwarfs the interval.} Per-snapshot
        host cycle: min \SI{2.02}{\second} / median \SI{9.75}{\second} / max
        \SI{23.57}{\second}. This is the binding cadence, $\sim$20$\times$ the
        \SI{500}{\milli\second} request at the median.
  \item \textbf{\texttt{pmemsave} dominates the cycle.} \texttt{pmemsave}
        duration: min \SI{0.76}{\second} / median \SI{7.06}{\second} / max
        \SI{21.57}{\second}; it is a median $81\%$ ($0.808$) of the host
        cycle. Throughput is $\approx$\SI{1.36}{\gibi\byte\per\second} on a
        quiet host, falling to
        \SIrange{0.27}{0.71}{\gibi\byte\per\second} under disk contention.
        Suspend ($\approx$\SI{0.35}{\second}) and resume
        ($\approx$\SI{0.05}{\second}) are negligible.
        % HONESTY (pmemsave-drift caveat): without a periodic flush /
        % delete-as-you-go cleanup, pmemsave DRIFTS within a run from
        % ~0.78 s to ~7.25 s as un-deleted 1 GiB dumps fill the imageDir and
        % page-cache writeback stalls each write. The drift is a disk-pressure
        % artifact, not an intrinsic pmemsave cost; the contended-throughput
        % range above reflects it. Before publication, separate the quiet-host
        % floor from the drift tail explicitly.
  \item \textbf{The VM is paused almost all the time, but not uniformly.}
        Estimated VM-pause fraction: min $0.74$ / median $0.946$ / max
        $0.978$.
        % HONESTY (pause-fraction spread): the pause fraction is NOT a single
        % constant -- it ranges 0.74-0.978 across ok cells. Report the spread,
        % not just the ~95% median, since pause fraction tracks pmemsave cost
        % and therefore inherits the drift caveat above.
  \item \textbf{Snapshots per cell scale with duration but are highly
        variable.} Per-duration (min/median/max):
        \SI{120}{\second} $\rightarrow$ 5 / 14 / 43;
        \SI{300}{\second} $\rightarrow$ 12 / 26 / 148;
        \SI{600}{\second} $\rightarrow$ 28 / 53 / 100.
        % HONESTY (7/20/28 minima correction): the often-quoted "7/20/28"
        % snapshot-count triplet is near the PER-DURATION MINIMA, not typical
        % values. The medians are 14/26/53. State this correction explicitly
        % in prose; do not cite 7/20/28 as representative.
  \item \textbf{Inertness, stated precisely.} Because the host snapshot cycle
        (median \SI{9.75}{\second}) dwarfs the \SI{500}{\milli\second}
        interval, and the guest already idles only $\approx$\SI{0.52}{\second}
        between snapshots, lowering \texttt{intervalMsec} below the
        snapshot-cycle floor cannot increase snapshot density -- the host
        cannot produce frames faster than \texttt{pmemsave} allows. The
        interval knob is therefore inert for throughput; it only sets the
        (already-honored) guest-time spacing. This motivates Paper B's
        single-interval design and reframes ``how often to sample'' as ``how
        cheaply can we snapshot.''
  \item \textbf{Validation by bug discovery (forward to the methods value).}
        \begin{itemize}
          \item Missing \texttt{bc}: the producer's post-resume sleep fell
                back to integer \texttt{sleep \$((intervalMsec/1000))} =
                \texttt{sleep 0} for any \texttt{intervalMsec} $< 1000$,
                yielding a measured guest gap of $\approx$\SI{6.6}{\milli\second}
                instead of the configured value. After the fix, measured
                guest gap $\approx$\texttt{intervalMsec} + \SI{25}{\milli\second}
                bash bookkeeping (e.g.\ \SI{125}{\milli\second} at
                \texttt{intervalMsec}=100).
                % HONESTY: 6.6 ms and 125 ms are from the small timing
                % experiments (3 and 19 snapshots), not the v3 campaign;
                % label them as instrumentation-run measurements, not
                % campaign statistics.
          \item Silent dump-cleanup failure: \texttt{unlink} failed under a
                root-owned \texttt{imageDir}, so every dump was retained and
                disk pressure inflated \texttt{pmemsave} (up to
                \SIrange{15}{31}{\second} on a near-full filesystem). Fixed
                with a \texttt{sudo}-fallback purge and producer self-clean.
          \item Methodological point: neither bug is visible without the six
                timestamps. The instrumentation is load-bearing precisely
                because it makes silent mis-calibration loud.
        \end{itemize}
  \item TODO: table -- \texttt{producer\_stats} aggregate (pmemsave, host
        cycle, suspend, resume, guest-run interval, pause fraction) with
        min/median/max columns over the 128 ok cells.
  \item TODO: figure -- per-snapshot \texttt{pmemsave\_sec} vs snapshot index
        for one quiet run and one drift-contaminated run, showing the
        disk-pressure tail.
  \item TODO: figure -- bar of mean cycle decomposition
        (suspend / pmemsave / flush / resume / guest-run sleep) showing
        pmemsave as the $\sim$81\% term and the \SI{500}{\milli\second}
        interval as a sliver.
\end{itemize}

%------------------------------------------------------------------
\section{Limitations}
\begin{itemize}
  \item Single requested interval in production. The v3 campaign fixed
        \texttt{intervalMsec} at \SI{500}{\milli\second}; the interval
        \emph{sweep} (100/250/500/1000/2000 ms) was a separate sub-pilot, and
        its analyzer-side metrics (F1, CV) were never produced because that
        sweep ran without \texttt{--keep-dumps}. The inertness claim rests on
        the cycle-vs-interval ratio, which is robust, but the per-family
        interval recommendations are design-justified, not analyzer-grounded.
        % HONESTY: the iv-recommendation table (iv_recommendations_v1.json) is
        % explicitly "design-justified, not empirically grounded by analyzer
        % output" -- C3 (n_windows >= 50) failed for every cell. Do not
        % present those recommendations as validated.
  \item Host-specific. All numbers are from a single capture host
        (\texttt{pcrserral}); \texttt{pmemsave} throughput and the suspend
        tail depend on the host's disk and libvirtd. The inertness direction
        generalizes (a full-RAM dump always dwarfs a sub-second interval), but
        the exact cycle distribution does not.
  \item Suspend latency tail. Under concurrent consumer load the
        instrumented ``suspend'' interval (which includes
        \texttt{domstate}-confirmation polling, not just the VCPU pause)
        spiked to tens of seconds in one run; the present instrumentation
        cannot split RPC vs polling vs actual pause. The campaign medians
        above are over producer-side runs, but the tail is real and
        under-characterized.
        % HONESTY: a single ~80 s suspend spike in one 19-snapshot run
        % dominates that run's mean; do not quote suspend means without the
        % median, and note the tail is sampled by too few events.
  \item Pmemsave drift conflates two regimes. The contended throughput range
        mixes the quiet-host floor with the disk-pressure drift; cleanly
        separating them needs a delete-as-you-go (or KEEP\_LAST=N) re-run.
\end{itemize}

%------------------------------------------------------------------
\section{Conclusion}
\begin{itemize}
  \item Per-snapshot timing instrumentation turned an opaque ``sampling
        interval'' knob into a measured, decomposed cost, and revealed that
        the requested guest interval is inert relative to the
        \texttt{pmemsave}-dominated host snapshot cycle.
  \item The binding cost is the hypervisor snapshot (\texttt{pmemsave}, a
        median $81\%$ of a median \SI{9.75}{\second} cycle), which both bounds
        temporal resolution and frees the interval knob. ``How often to
        sample'' is therefore the wrong question; ``how cheaply can we
        snapshot'' is the right one.
  \item This methods finding is the foundation the rest of the portfolio
        stands on: Paper A's cost model, Paper B's single-interval (gated)
        design, Paper F's streaming capture system, and the Plan 05
        snapshot-throughput proposal all inherit the inertness result
        established here.
\end{itemize}

\bibliographystyle{IEEEtran}
\bibliography{../refs}
\end{document}
